\documentclass[11pt]{article}
\usepackage{preamble}

\begin{document}

% ======================================================================
% Title page
% ======================================================================
\begin{titlepage}
\centering
\vspace*{1.5cm}
{\Large ARTEMIS System Design}\\[0.4cm]
{\LARGE\bfseries System Architecture}\\[0.3cm]
{\large Revision 2}\\[0.6cm]
{\large FAST NUCES Final-Year Project \textbf{F26-008}}\\[0.2cm]
{\small \textit{Adaptive Reasoning, Task Execution and Multi-agent Intelligence System}}\\[1.1cm]

\begin{tabular}{@{}ll@{}}
\toprule
Team member & Roll number \\
\midrule
Hashir Rauf & 23L-2572 \\
Hira Khalid & 23L-2594 \\
Doureesha Batool & 23L-2651 \\
\bottomrule
\end{tabular}

\vspace{0.7cm}

{\normalsize Supervisor: Mr.\ Mubasher Hussain}

\vfill
{\normalsize Internal design document. Not for external distribution.}\\[0.3cm]
{\normalsize Compiled \today}
\end{titlepage}

% ======================================================================
% Front matter
% ======================================================================
\pagenumbering{roman}
\setcounter{page}{1}
\tableofcontents
\clearpage
\pagenumbering{arabic}
\setcounter{page}{1}

% ======================================================================
\section{Purpose and Status of This Document}
% ======================================================================

This document specifies the proof-of-concept architecture for ARTEMIS. It supersedes the
outline architecture recorded in the prototype development plan, which is referred to
throughout as Revision 1.

Revision 1 described the system as six stacked boxes with an AI gateway at the centre.
That description was adequate for communicating intent but inadequate as an engineering
specification, for three reasons. The gateway was a single component responsible for
routing, context assembly, memory, planning, policy and dispatch simultaneously, which is
neither testable nor separable. Security appeared as a labelled box at the foot of the
diagram rather than as a property of the whole system. And the remote-access facility was
drawn inside the core stack, which weakened the claim that ARTEMIS works without a
network.

Revision 2 addresses all three. It decomposes the gateway into four cooperating
components behind a single execution chokepoint, promotes security and performance to
cross-cutting planes, and draws the remote-access facility outside the core so that its
removal provably changes nothing about local operation. It also introduces one component
Revision 1 did not name at all, the Policy Engine, which is described in
Section~\ref{sec:policy} and which is the mechanism by which this project's central
research finding becomes enforced structure rather than stated intention.

Section~\ref{sec:changes} lists every change from Revision 1 with its rationale.

% ======================================================================
\section{Architectural Thesis}
% ======================================================================

Every structural decision in this document follows from a single sentence.

\begin{finding}[Design thesis]
ARTEMIS is local-first, cloud-optional and remotely accessible, but never
cloud-dependent for core operation, and never in possession of scope the user has not
explicitly granted.
\end{finding}

This is not a deployment preference. It is the architectural encoding of the project's
central research hypothesis: that restricting an assistant's scope increases user trust
and adoption without meaningfully reducing its usefulness. If that hypothesis is to be
tested honestly, the restriction has to be structural rather than behavioural. It must be
impossible by construction, not merely discouraged by instruction. A boundary that the
system could cross but chooses not to is a policy; a boundary the system has no code path
across is an architecture. This document is concerned throughout with the second kind.

\subsection{Degradation contract}

The correctness of a local-first system is defined by how it behaves as capabilities are
withdrawn from it. The rows of Table~\ref{tab:degradation} are acceptance criteria, each
one testable, not aspirations.

\begin{table}[htbp]
\centering
\small
\caption{Required behaviour under withdrawal of each external capability}
\label{tab:degradation}
\begin{tabular}{@{}p{0.30\textwidth}p{0.62\textwidth}@{}}
\toprule
\textbf{Withdrawn capability} & \textbf{Required behaviour} \\
\midrule
Internet connectivity lost &
Full local operation continues. Only cloud model calls, web search and mail sending are
disabled, and each fails as an explicit, retryable queued action rather than a silent
drop. \\
\addlinespace
All cloud model APIs unavailable &
Falls back to the configured local model. Features exceeding local capability degrade to
a stated limitation, not an error. \\
\addlinespace
Local model unavailable &
Falls forward to the cloud chain, if and only if the user has enabled cloud fallback for
that workspace. \\
\addlinespace
Remote device disconnects &
The host continues uninterrupted. Approvals in flight revert to pending and surface on
the host interface. \\
\addlinespace
Cloudflare edge unreachable &
Local operation is entirely unaffected, because the remote plane is not on the execution
path. \\
\addlinespace
Workspace not explicitly granted &
ARTEMIS cannot see it. This is not a refusal to act; the path is not resolvable through
the Workspace Broker at all. \\
\bottomrule
\end{tabular}
\end{table}

\subsection{Non-goals}

An architecture is defined as much by what it refuses as by what it provides. The
following are permanently out of scope and are enforced, not merely avoided.

\begin{itemize}
  \item \textbf{No hosted user data.} Documents, embeddings, session history, memory,
        credentials and audit logs never leave the host machine. The only outbound
        telemetry is anonymous, opt-in and schema-fixed diagnostic logging.
  \item \textbf{No autonomous outward action.} No message is sent, no web request issued
        and no file modified without passing an approval gate.
  \item \textbf{No ambient surveillance.} There is no background browsing, screen
        capture, keystroke logging or always-listening recording. A voice session is
        explicitly armed by the user and visibly indicated while active.
  \item \textbf{No self-modifying execution.} The agent cannot write and then run
        arbitrary code against the host as a step in its own plan.
\end{itemize}

% ======================================================================
\section{Layered Overview}
% ======================================================================

The system is organised as six layers on the execution path, with two planes crossing all
of them. Figure~\ref{fig:layers} shows the whole arrangement.

\begin{figure}[htbp]
\centering
\includegraphics[width=0.72\textwidth]{fig_layers.pdf}
\caption{The ARTEMIS layered architecture. Layers 1 to 5 form the local execution path.
The security and performance planes are drawn as enclosing bands rather than as layers
because they are properties every layer must hold. Layer 6 is drawn separately, in
Figure~\ref{fig:remote}, because nothing in Layers 1 to 5 depends on it.}
\label{fig:layers}
\end{figure}

\subsection{Why two planes rather than two further layers}

Security and performance are not stages through which a request passes. They are
properties that every stage must hold simultaneously. Drawing them as layers would imply
that a request can be past security, which is precisely the reasoning error that produces
confused-deputy defects: a component trusts an input because some earlier component is
assumed to have checked it. Drawing them as planes states the opposite, that every arrow
in the figure crosses both, and that no component may assume a check has already
happened.

\subsection{Component inventory}

Table~\ref{tab:components} lists all eighteen components. The final column records the
single constraint each component exists to hold, which is the property whose violation
would break the design thesis of Section~2.

\begin{table}[htbp]
\centering
\small
\caption{Component inventory, with the invariant each component holds}
\label{tab:components}
\begin{tabular}{@{}rp{0.23\textwidth}cp{0.26\textwidth}p{0.26\textwidth}@{}}
\toprule
\textbf{\#} & \textbf{Component} & \textbf{Layer} & \textbf{Responsibility} & \textbf{Must never} \\
\midrule
1 & Command-line interface & 1 & Scriptable control surface & Bypass the approval broker \\
2 & Graphical interface & 1 & Visual control panel & Bind to a non-loopback interface \\
3 & Control panel view model & 1 & Single source of interface truth & Hold state the backend does not own \\
4 & Voice Session Manager & 2 & Detection, transcription, speech output & Record when not explicitly armed \\
5 & Text Stream Manager & 2 & Token streaming and cancellation & Buffer past a cancellation \\
6 & Turn Normaliser & 2 & Modality-neutral turn object & Leak modality downstream \\
7 & Model Router & 3 & Model selection and failover & Send workspace data to a disabled tier \\
8 & Context Engine & 3 & Bounded, cited context assembly & Include content outside the workspace \\
9 & Policy Engine & 3 & Risk classification and gating & Be consulted after execution \\
10 & Agent Planner & 3 & Multi-step plan construction & Execute anything itself \\
11 & Tool Dispatcher & 3 & Sole execution chokepoint & Accept a call lacking a verdict \\
12 & Capability services & 4 & The nine product features & Touch the filesystem directly \\
13 & Workspace Broker & 4 & Path resolution and confinement & Resolve a path outside a grant \\
14 & Indexer & 4 & Incremental embedding and indexing & Block the interactive path \\
15 & Undo Journal & 4 & Snapshots and inverse operations & Record an action after the fact \\
16 & Approval Broker & 4 & Pending actions across surfaces & Auto-approve on timeout \\
17 & Data plane stores & 5 & Local persistence & Serialise outside the host \\
18 & Remote access plane & 6 & Authenticated relay & Store or read plaintext user data \\
\bottomrule
\end{tabular}
\end{table}

% ======================================================================
\section{The Request Lifecycle}
% ======================================================================

Figure~\ref{fig:lifecycle} shows the single most important path in the system. Every
request, whether spoken or typed, and whether raised locally or from a paired remote
device, traverses exactly this sequence. There are no shortcuts and no privileged
entry points.

\begin{figure}[htbp]
\centering
\includegraphics[width=0.86\textwidth]{fig_lifecycle.pdf}
\caption{The canonical request lifecycle. Modality is erased at the Turn Normaliser, so
every stage below it is identical for voice and for text. The gated branch returns to the
trunk beneath the Undo Journal, so that an approved action is snapshotted before it is
dispatched.}
\label{fig:lifecycle}
\end{figure}

Five invariants hold along this path, and each is directly testable.

\begin{enumerate}
  \item The Policy Engine is consulted before the Tool Dispatcher, never after.
  \item The Undo Journal snapshots before mutation, never after.
  \item An audit entry is written even when the action is rejected, so that the log
        records what was proposed as well as what was performed.
  \item A tool call carrying no policy verdict is rejected by the dispatcher.
  \item Voice and text turns are indistinguishable from the planner onward.
\end{enumerate}

The fifth invariant is the mechanism behind the voice-parity claim discussed in
Section~\ref{sec:services}, and it is worth stating plainly why it matters. If the voice
path were a separate implementation, then every approval gate would exist twice, and the
two copies would drift. Because modality is erased before the planner, the gates exist
once and voice inherits them without further work.

% ======================================================================
\section{Layer 3: The AI Gateway}
% ======================================================================

Revision 1 assigned routing, context assembly, memory, planning, policy and dispatch to a
single gateway component. Revision 2 separates these into four components that cooperate
through one shared execution chokepoint, as shown in Figure~\ref{fig:gateway}.

\begin{figure}[htbp]
\centering
\includegraphics[width=0.68\textwidth]{fig_gateway.pdf}
\caption{Internal structure of the AI gateway. The planner produces a plan but never
executes it; the Tool Dispatcher executes but never decides. Separating the two is what
makes the pause before any workspace-changing call structural rather than conventional.}
\label{fig:gateway}
\end{figure}

\subsection{Model Router}

The Model Router owns tier selection and, more importantly, data-egress enforcement. It
is the only component that knows a given workspace's cloud policy, which makes it the
correct place to block egress. Placing that check in the tool layer would be too late,
and placing it in the prompt would not be a check at all.

\begin{figure}[htbp]
\centering
\begin{minipage}[t]{0.47\textwidth}
\centering
\includegraphics[width=\linewidth]{fig_router.pdf}
\caption{Tier selection and the egress gate. Tiers above zero are reachable only when the
active workspace permits cloud use; otherwise the router degrades and explains rather
than silently succeeding.}
\label{fig:router}
\end{minipage}\hfill
\begin{minipage}[t]{0.47\textwidth}
\centering
\includegraphics[width=\linewidth]{fig_context.pdf}
\caption{Context assembly. Every retrieved chunk carries the identifier of the workspace
it came from; any chunk not matching the active grant is dropped before budgeting, and
the drop is recorded.}
\label{fig:context}
\end{minipage}
\end{figure}

The fallback order from the proposal is retained: local models first, then Gemini, then
OpenAI, then Anthropic. Any endpoint offering an OpenAI-compatible interface can be
registered as a further tier through configuration alone, so adding a provider requires
no code change. Failover is bounded at one retry per tier before advancing, so a
workspace that has disabled cloud access has exactly one tier and fails cleanly rather
than hanging through a chain of unreachable providers.

\subsection{Context Engine}

The Context Engine assembles the working context for a request and enforces the bounded
scope hypothesis on the read side. The Workspace Broker, described in
Section~\ref{sec:broker}, enforces the same boundary on the write side.


The citation map is what makes citation-grounded synthesis verifiable rather than merely
plausible. It is carried through to the rendered answer, so a claim that cannot be traced
back to a specific chunk, file and location is marked as unsupported rather than
presented alongside supported claims. This directly addresses the hallucination and
citation-accuracy themes of the literature review.

\subsection{Policy Engine}
\label{sec:policy}

The Policy Engine is the most significant addition in this revision. It converts the
research finding that users want approval-first behaviour, configurable to their own
tolerance, from an interface convention into an enforced architectural component.

Every proposed tool call is classified into exactly one of four tiers.

\begin{figure}[htbp]
\centering
\begin{minipage}[t]{0.47\textwidth}
\centering
\includegraphics[width=\linewidth]{fig_tiers.pdf}
\caption{The four risk tiers. Only the reversible tier is user-configurable; outward
actions are always gated, and destructive actions are refused structurally.}
\label{fig:tiers}
\end{minipage}\hfill
\begin{minipage}[t]{0.47\textwidth}
\centering
\includegraphics[width=\linewidth]{fig_policy_order.pdf}
\caption{Policy evaluation order. Denials are evaluated before gates, and gates before
permissions, so that a refusal cannot be reached around.}
\label{fig:policyorder}
\end{minipage}
\end{figure}

\begin{finding}[Why the destructive tier has no dialogue]
There is no approval dialogue for tier three because there is no code path to it. This is
what makes the claim that ARTEMIS deletes nothing a property of the system rather than a
promise about the model's behaviour. A gate that can be argued past is a policy; a gate
that does not exist is an architecture.
\end{finding}

Classification proceeds in a fixed order, shown in Figure~\ref{fig:policyorder}, in which
the first matching rule wins. Ordering the denials first means that no later rule can
grant what an earlier rule refused.


Figure~\ref{fig:approval} works the classification through end to end for a
representative request.

\begin{figure}[htbp]
\centering
\includegraphics[width=0.62\textwidth]{fig_approval.pdf}
\caption{A worked classification and approval. Read operations execute automatically;
the two reversible operations are presented together, at the granularity the user
actually asked about, with the undo window and the absence of egress stated explicitly.}
\label{fig:approval}
\end{figure}

The option to approve and remember is how the configurable half of the research finding
is implemented. Selecting it writes a scoped rule covering this operation type, in this
workspace, for this shape of action. It does not lower the user's protection globally,
which is the failure mode that makes permission systems degrade into permanent consent.

\subsection{Agent Planner}

The planner is built on LangGraph, as recorded in the project's stack decisions. The
architectural constraint on it is simple and absolute: it emits a plan and never executes
one. The plan is data, an ordered list of proposed tool calls, and remains inert until
the dispatcher acts on it.

The reflection loop, by which the planner revises its approach after observing a result,
is bounded by an explicit cycle cap and a wall-clock budget. A planner that cannot
converge within those bounds surfaces its partial plan and asks the user, rather than
continuing to iterate.

% ======================================================================
\section{Layer 4: Capability Services}
\label{sec:services}
% ======================================================================

The nine features from the project scope map onto capability services. Each is a plugin
behind a uniform tool contract, and none touches the filesystem, the network or the
database directly. All access passes through the shared substrate shown in
Figure~\ref{fig:services}.

\begin{figure}[htbp]
\centering
\begin{minipage}[t]{0.47\textwidth}
\centering
\includegraphics[width=\linewidth]{fig_services.pdf}
\caption{Capability services are reached only through the Tool Dispatcher, and reach
storage only through the shared substrate. Each service is coloured by the highest risk
tier it may request.}
\label{fig:services}
\end{minipage}\hfill
\begin{minipage}[t]{0.47\textwidth}
\centering
\includegraphics[width=\linewidth]{fig_broker.pdf}
\caption{Path resolution and confinement. Canonicalisation occurs before the bounds
assertion, which is what prevents symbolic-link and relative-traversal escapes.}
\label{fig:broker}
\end{minipage}
\end{figure}

\begin{table}[htbp]
\centering
\small
\caption{The nine capability services and the highest risk tier each may request}
\label{tab:services}
\begin{tabular}{@{}clcp{0.50\textwidth}@{}}
\toprule
\textbf{\#} & \textbf{Service} & \textbf{Max tier} & \textbf{Constraint} \\
\midrule
F1 & Session Resume & T0 & No model involvement in the first version; a deterministic replay of recorded session state. \\
F2 & Folder Tidy & T1 & Moves and groups only. It cannot delete, because no code path exists to do so. \\
F3 & Notes Synthesis & T0 & Reads and generates. The output remains a proposal until the drafting service saves it. \\
F4 & Pattern Watcher & T1 & Detects three or more repetitions and may only propose an automation rule. It cannot install one. \\
F5 & Code Assistant & T0 & Explains, and produces a diff. It never writes to a file and never executes code. \\
F6 & Web Search & T2 & On demand only. There is no background browsing, and every query is logged. \\
F7 & Mail Composer & T2 & Draft, approve, then send. Sending is gated regardless of user configuration. \\
F8 & Document Drafter & T1 & Always writes a new file. It never overwrites an existing one. \\
F9 & Speech Input and Output & n/a & Not a capability service at all, but a Layer 2 modality over identical gates. \\
\bottomrule
\end{tabular}
\end{table}

\begin{finding}[Voice parity is structural]
Speech is deliberately not implemented as a capability service. It is normalised at
Layer 2 into the same turn object as typed input, and therefore inherits every approval
gate automatically. Parity with the text interface is achieved by construction rather
than by reimplementing the gates along a second path, which addresses the voice
interaction parity theme of the literature review.
\end{finding}

\subsection{Workspace Broker}
\label{sec:broker}

The Workspace Broker is the only component that converts a logical reference into a real
filesystem path. It is therefore the enforcement point for bounded scope on the write
side, and its correctness is load-bearing for the entire trust claim.


The ordering of steps two and three is not incidental. Resolving symbolic links and
relative segments before asserting containment is the difference between a boundary and a
suggestion; performing the assertion first would allow a path that appears to be inside
the grant to resolve to a location outside it.

\subsection{Undo Journal}

Before any reversible operation, the Undo Journal records both a forward and an inverse
operation, choosing a snapshot strategy according to cost: full content snapshots for
small files, copy-on-write links for large ones, and inverse-operation records alone for
moves and renames, which need no content copy.

Undo operates at the granularity of a plan rather than of an individual operation.
Reversing a folder tidy restores all twelve moves as a single unit and in reverse order,
because a plan is the unit the user was asked to approve. Per-operation undo remains
available but is secondary, since it does not correspond to anything the user consented
to as a whole.

% ======================================================================
\section{Layer 5: The Data Plane}
% ======================================================================

All persistent state is local. The store layout separates the relational system of
record, the per-workspace vector indexes, the content-addressed snapshot blobs, a
discardable cache, and non-secret configuration. Secrets are held in the operating
system keyring and never written into the ARTEMIS directory tree.

Per-workspace index isolation is a deliberate choice rather than an organisational
convenience. Because each workspace has its own index, the embeddings of one workspace
cannot be retrieved while another is active, even in the presence of a query defect: the
index is simply not open. It also reduces workspace revocation to a directory deletion,
and it bounds the index rebuild cost that the prototype plan identified as a feasibility
risk.

\subsection{Memory architecture}

The memory design is the clearest differentiator of the system and is shown in
Figure~\ref{fig:memory}.

\begin{figure}[htbp]
\centering
\begin{minipage}[t]{0.47\textwidth}
\centering
\includegraphics[width=\linewidth]{fig_memory.pdf}
\caption{Five kinds of memory, distinguished by lifetime and by store, feeding a single
context engine that merges, budgets and bounds-asserts before anything reaches a model.}
\label{fig:memory}
\end{minipage}\hfill
\begin{minipage}[t]{0.47\textwidth}
\centering
\includegraphics[width=\linewidth]{fig_memory_vs_index.pdf}
\caption{The file index and ARTEMIS memory answer different questions and have different
natures. Conflating them is the design error that turns a bounded assistant into a
surveillance tool.}
\label{fig:memvsidx}
\end{minipage}
\end{figure}

A distinction that the design depends on, and which the final report should state
explicitly, is that between the file index and memory proper.


The file index is mechanical. It records what exists and where, is derived from the
filesystem, holds paths and content hashes rather than interpretations, and can be
rebuilt at any time without loss. Memory is interpretive. It records what the system has
learned about a workspace, a task and a user's preferences, accumulates over sessions,
and is genuinely lost if deleted. The first requires only a grant; the second requires
ongoing consent, and every entry in it is individually visible and individually
deletable.

\subsection{The disclosure panel}

The panel that shows what ARTEMIS currently holds is not a summary the system writes
about itself. It is a live query across every store: granted workspaces, indexed files,
embedded chunks, remembered facts, retained sessions, automation rules, policy exceptions,
actions taken with their undo state, every record of data sent to a cloud provider, and
the names of held credentials.

Every row is deletable from the panel, and deletion cascades: removing a workspace
removes its index directory, its file rows, its embeddings and its memory entries. This
matters because it makes the trust claim falsifiable. A user can check it rather than
believe it, which is a stronger position than any assurance the interface could offer.

% ======================================================================
\section{Layer 6: The Remote Access Plane}
% ======================================================================

The remote access plane is drawn outside the core stack because nothing in Layers 1 to 5
depends on it. Removing it entirely leaves a fully functional single-machine system, and
that is the test of whether the separation is real rather than notional.

\begin{figure}[htbp]
\centering
\includegraphics[width=0.48\textwidth]{fig_remote.pdf}
\caption{The remote access plane. The durable object records which authenticated device
is connected to which host and relays encrypted frames, but holds no plaintext, no
workspace content and no keys.}
\label{fig:remote}
\end{figure}

\subsection{Pairing}

Pairing establishes an end-to-end encrypted channel that the relaying edge cannot read,
and requires a confirmation performed physically at the host.

\begin{figure}[htbp]
\centering
\begin{minipage}[t]{0.47\textwidth}
\centering
\includegraphics[width=\linewidth]{fig_pairing.pdf}
\caption{The pairing handshake. The host-side confirmation is deliberate: possession of
the pairing code alone is never sufficient, so compromising the edge cannot pair a
device.}
\label{fig:pairing}
\end{minipage}\hfill
\begin{minipage}[t]{0.47\textwidth}
\centering
\includegraphics[width=\linewidth]{fig_topology.pdf}
\caption{Process topology on the host machine. The core process is the only holder of
user data; the interfaces are clients, and the remote bridge dials outward.}
\label{fig:topology}
\end{minipage}
\end{figure}

\subsection{Constraints on remote sessions}

The remote plane is a control surface, not a privilege escalation. The following
constraints are enforced at the host, not at the edge, because the edge is not trusted.

\begin{itemize}
  \item A remote session runs at privilege equal to or lower than a local session, never
        higher.
  \item A workspace may be marked local-only, in which case it is invisible to remote
        sessions entirely.
  \item Outward actions may be configured to require approval at the host regardless of
        where they were requested. This is enabled by default.
  \item Sessions expire on a fixed lifetime and on host sleep. There are no indefinite
        sessions.
  \item Every remote-originated action is tagged with its originating device in the audit
        log.
  \item Revocation is immediate and unilateral from the host: the session is terminated
        and keys are rotated.
  \item If the edge is unreachable the host simply has no remote clients, and nothing
        else about its operation changes.
\end{itemize}

% ======================================================================
\section{The Security and Trust Plane}
% ======================================================================

Table~\ref{tab:security} lists every security control with the component that enforces it
and the literature review theme that motivated it. A control with no named enforcement
point is an intention rather than a mechanism, which is why the middle column exists.

\begin{table}[htbp]
\centering
\small
\caption{Security controls, their enforcement points, and their research motivation}
\label{tab:security}
\begin{tabular}{@{}p{0.30\textwidth}p{0.34\textwidth}p{0.28\textwidth}@{}}
\toprule
\textbf{Control} & \textbf{Enforced at} & \textbf{Research link} \\
\midrule
Workspace confinement & Workspace Broker & Bounded agents and trust \\
Explicit grant model & Workspace records and interface & Bounded agents and trust \\
Tool allowlisting & Tool Dispatcher & Agentic attack surface \\
Approval gates & Policy Engine & Planning and undo \\
Risk classification & Policy Engine & Bounded agents and trust \\
Destructive-operation refusal & Policy tier three, no code path & Bounded agents and trust \\
Encryption at rest & Encrypted database and keyring & Bounded agents and trust \\
Encryption in transit & Transport security and end-to-end relay & Bounded agents and trust \\
Credential vault & Operating system keyring & Bounded agents and trust \\
Prompt injection defence & Layered, see Section~\ref{sec:injection} & Agentic attack surface \\
Egress control & Model Router egress gate & Privacy and local-first operation \\
Append-only audit log & Hash-chained log & Assurance frameworks \\
Single-action undo & Undo Journal & Planning and undo \\
Remote device authentication & Worker with host confirmation & Bounded agents and trust \\
Session expiry & Session lifetime and sleep hook & Bounded agents and trust \\
Least privilege & Remote no greater than local & Bounded agents and trust \\
Citation grounding & Context Engine citation map & Hallucination and citation accuracy \\
Human review of code changes & Code assistant emits diffs only & Code assistance with review \\
\bottomrule
\end{tabular}
\end{table}

\subsection{Prompt injection}
\label{sec:injection}

ARTEMIS reads documents, source code and web results, all of which can be influenced by
an adversary. The defence is not detection, which is unreliable, but a structural
guarantee that retrieved content can never become an authorised instruction.

\begin{figure}[htbp]
\centering
\includegraphics[width=0.66\textwidth]{fig_injection.pdf}
\caption{Five independent layers of prompt injection defence. Each is sufficient to
prevent a class of attack, and the final layer bounds the damage of any injection that
somehow defeats the preceding four.}
\label{fig:injection}
\end{figure}

\begin{finding}[The strongest form of the hypothesis]
The fifth defence layer is what makes bounded scope load-bearing rather than rhetorical.
Confining the agent to one granted workspace bounds the worst case, not merely the
expected case. That is an argument no amount of prompt engineering can make, and it is
the clearest evidence this project can offer for its central hypothesis.
\end{finding}

\subsection{Threat model}

\begin{table}[htbp]
\centering
\small
\caption{Threat model, with mitigations and honestly stated residual risk}
\label{tab:threats}
\begin{tabular}{@{}p{0.20\textwidth}p{0.20\textwidth}p{0.32\textwidth}p{0.20\textwidth}@{}}
\toprule
\textbf{Adversary} & \textbf{Capability assumed} & \textbf{Mitigation} & \textbf{Residual risk} \\
\midrule
Malicious document or web page & Injects instructions into context & The five layers of Section~\ref{sec:injection} & The user approves a plausible but harmful action \\
\addlinespace
Compromised model provider & Observes what is sent & Local-first default, per-workspace opt-in, redaction, egress logging & Content sent under explicit opt-in \\
\addlinespace
Compromised edge & Observes relayed traffic and controls the worker & End-to-end encryption, host-side pairing confirmation, no plaintext at the relay & Traffic metadata such as timing and volume \\
\addlinespace
Stolen paired device & Holds a valid session & Session expiry, host revocation, local-only workspaces, host approval for outward actions & Actions taken before revocation \\
\addlinespace
Local attacker with disk access & Reads the ARTEMIS directory & Encryption at rest, keyring for secrets & Attack against a running process's memory \\
\addlinespace
The project team & Operates the telemetry endpoint & No user data is collected by design; logs are schema-fixed and opt-in & None for user content \\
\bottomrule
\end{tabular}
\end{table}

% ======================================================================
\section{The Execution and Performance Plane}
% ======================================================================

The governing rule is that the interactive path stays thin and everything expensive runs
asynchronously behind it.

\begin{figure}[htbp]
\centering
\includegraphics[width=0.66\textwidth]{fig_performance.pdf}
\caption{The interactive path never blocks on background work. Background progress is
always visible, pausable and cancellable, so that waiting is legible rather than
mysterious.}
\label{fig:performance}
\end{figure}

Work is scheduled at five priorities: user-initiated execution, which preempts everything;
retrieval for an active turn, which is immediate but budget-capped; incremental index
updates for changed files, debounced by roughly two seconds; initial workspace ingestion,
which is batched, throttled and resumable; and maintenance, which runs only when idle.

Several strategies follow from this structure. Indexing is incremental, so that editing
one line re-embeds one chunk rather than an entire file. Retrieval is two-stage, pairing
cheap vector recall with a reranking pass, which keeps the index small while preserving
quality. The local model is kept warm with an idle-unload timer so that first-token
latency does not include a cold load. Every queue is bounded and applies backpressure, so
that a saturated system sheds the lowest-priority work visibly rather than falling
silently behind. Resource ceilings on memory, index size and processor share are
configurable, because a background assistant on a student laptop must never be the reason
the machine feels slow.

% ======================================================================
\section{Process and Deployment Topology}
% ======================================================================


Three properties of this arrangement are deliberate. The core process is the only holder
of user data, so terminating an interface loses nothing. There is no inbound listening
port beyond loopback, because the remote bridge establishes an outbound connection, which
means there is no port to scan, no traversal to negotiate and no firewall exception to
request. And the local model server runs as a separate process, so that a model crash
cannot take down the workspace layer and models can be replaced without touching ARTEMIS
itself.

% ======================================================================
\section{Traceability Between Architecture and Research}
% ======================================================================

Table~\ref{tab:traceability} records why each non-obvious structural decision exists. It
is the defensible core of the architecture chapter, because it demonstrates that the
design follows from the research rather than being justified after the fact.

\begin{table}[htbp]
\centering
\small
\caption{Each structural decision and the finding that produced it}
\label{tab:traceability}
\begin{tabular}{@{}p{0.42\textwidth}p{0.44\textwidth}@{}}
\toprule
\textbf{Architectural decision} & \textbf{Driven by} \\
\midrule
Workspace Broker as the sole path resolver & The bounded scope hypothesis, and questionnaire results on workspace trust \\
\addlinespace
Policy Engine as a distinct component & The approval-first, configurable finding \\
\addlinespace
Destructive tier with no code path & The requirement that deleting nothing be a property rather than a promise \\
\addlinespace
Undo at plan granularity & Recovery is meaningful at the unit the user approved \\
\addlinespace
Voice normalised at Layer 2 & Voice interaction parity, achieved by construction \\
\addlinespace
Citation map through the Context Engine & Hallucination and citation accuracy themes \\
\addlinespace
Code assistant emits diffs only & Code assistance requires human review \\
\addlinespace
Local-first with per-workspace egress opt-in & The local versus cloud split in questionnaire responses \\
\addlinespace
Provenance tagging of context & The agentic attack surface literature \\
\addlinespace
Session resume implemented without a model & Task interruption and resumption; also the lowest-risk first feature \\
\addlinespace
Durable object holds no plaintext & The privacy claim must survive compromise of the edge \\
\addlinespace
Remote privilege no greater than local & Remote convenience must not widen the trust boundary \\
\addlinespace
Background queue with visible progress & Trust requires legibility of what the system is doing \\
\bottomrule
\end{tabular}
\end{table}

% ======================================================================
\section{Build Sequence}
% ======================================================================

The staged build sequence from the prototype plan is retained, adjusted so that the two
components this revision adds are built early. Both are load-bearing, and both are far
cheaper to build before there are features to retrofit them into.

\begin{table}[htbp]
\centering
\small
\caption{Revised build sequence}
\label{tab:build}
\begin{tabular}{@{}cp{0.40\textwidth}p{0.42\textwidth}@{}}
\toprule
\textbf{Stage} & \textbf{Deliverable} & \textbf{Why at this point} \\
\midrule
1 & Workspace foundation: database, Workspace Broker, file watcher and disclosure panel, with no model involvement & The boundary must exist before anything can be bounded \\
\addlinespace
2 & Policy Engine, Approval Broker and Undo Journal, exercised with placeholder tools & Cheap now and very expensive to retrofit; provable before any model is introduced \\
\addlinespace
3 & AI gateway: Model Router with the local tier and one cloud tier, Context Engine, Agent Planner and Tool Dispatcher & The highest-risk stage, validated against gates that already work \\
\addlinespace
4 & Core file features: session resume, folder tidy and pattern watching & The first real capabilities, exercising the read and reversible tiers end to end \\
\addlinespace
5 & Understanding features: notes synthesis and code assistance, with the index, retrieval and citation map & Adds the vector index once the boundary is proven \\
\addlinespace
6 & Outward features: web search, mail and document drafting & The first outward-tier actions, introduced only after the gates are demonstrated \\
\addlinespace
7 & Speech input and output over the existing Turn Normaliser & Should require no new gates. If it does, Layer 2 was designed wrongly, which makes this stage the test of the parity claim \\
\addlinespace
8 & Remote access plane & Genuinely optional; a complete product exists without it \\
\bottomrule
\end{tabular}
\end{table}

\begin{finding}[The most consequential scheduling change]
Stage two now precedes all model work. This means the trust machinery is testable, and
demonstrable to a supervisor, before any model non-determinism enters the system. A
failing approval gate at that stage is a deterministic bug with a short reproduction,
which is not true of the same defect discovered after the planner exists.
\end{finding}

% ======================================================================
\section{Open Questions}
% ======================================================================

The prototype plan recorded four feasibility questions, which remain open. Each is
restated here with the architectural decision it blocks and the stage by which it must be
settled.

\begin{table}[htbp]
\centering
\small
\caption{Open questions carried forward, with what each one blocks}
\label{tab:open}
\begin{tabular}{@{}p{0.36\textwidth}p{0.40\textwidth}c@{}}
\toprule
\textbf{Question} & \textbf{Blocks} & \textbf{By stage} \\
\midrule
Which local model size is realistic on a typical student laptop? & The default local tier, and the resource envelope & 3 \\
\addlinespace
Which desktop interface approach is most reliable across features? & Layer 1. The current inclination is the local webview, since one view model can then serve both the local and remote interfaces & 1 \\
\addlinespace
Are local speech models accurate enough to avoid cloud fallback? & Whether voice input ever becomes an outward-tier action & 7 \\
\addlinespace
How large can a workspace index grow before retrieval or rebuild degrades? & Per-workspace index sizing and the two-stage retrieval thresholds & 5 \\
\bottomrule
\end{tabular}
\end{table}

This revision raises three further questions. Whether encryption at rest should use an
encrypted database or rely on operating-system disk encryption together with the keyring,
which determines how the encryption claim in Table~\ref{tab:security} may be stated.
Whether a reranking model is affordable locally, or whether vector-only retrieval is
sufficient, which determines the two-stage retrieval design. And whether the relay needs
message durability at all, since a live-only channel is both simpler and a stronger form
of the claim that the relay holds no data.

% ======================================================================
\section{Summary of Changes from Revision 1}
\label{sec:changes}
% ======================================================================

\begin{table}[htbp]
\centering
\small
\caption{Changes from Revision 1, with rationale}
\label{tab:changes}
\begin{tabular}{@{}cp{0.42\textwidth}p{0.42\textwidth}@{}}
\toprule
\textbf{\#} & \textbf{Change} & \textbf{Rationale} \\
\midrule
1 & The AI gateway is decomposed into router, context, policy and planner components behind one dispatcher & The single-component design was untestable and difficult to defend \\
2 & A Policy Engine is introduced as a first-class component with four risk tiers & Converts the approval-first finding into enforced structure \\
3 & Security is promoted from a component to a cross-cutting plane & Security is a property of every interaction, not a stage \\
4 & Execution and performance is added as a second cross-cutting plane & Makes the synchronous and asynchronous split explicit and testable \\
5 & The remote plane is drawn outside the core stack & Demonstrates the privacy claim: removing it changes nothing locally \\
6 & Memory is formally separated from the file index & The clearest differentiator, and the clearest privacy boundary \\
7 & A Turn Normaliser is introduced at Layer 2 & Makes voice parity structural rather than duplicated \\
8 & The Workspace Broker becomes the sole path resolver & Turns bounded scope from a policy into an invariant \\
9 & Undo moves to plan granularity & Matches the unit the user was asked to approve \\
10 & Prompt injection is addressed as an explicit five-layer defence & Answers the central threat identified in the review \\
11 & The build sequence is reordered to place trust machinery before model work & Cheap now, nearly impossible to retrofit later \\
12 & A traceability matrix is added & Ties every structural choice to a cited finding \\
\bottomrule
\end{tabular}
\end{table}

\end{document}
